\subsection{Purpose}

Change-management procedures ensure that changes to systems, applications, networks, devices, configurations, data, and technology services are authorized, evaluated, tested, implemented, and documented in a controlled manner.

The purpose of change management is to reduce service interruptions, security weaknesses, data loss, configuration errors, and unintended effects on other systems.

\subsection{Change Requirements}

Changes \must{} be managed when they may affect:
\begin{itemize}
    \item System security or access controls
    \item System availability or performance
    \item Network connectivity or routing
    \item Applications, databases, or stored data
    \item Operating systems, firmware, or installed software
    \item Backup, logging, monitoring, or recovery capabilities
    \item Encryption, authentication, or authorization
    \item Compliance, contractual, or legal requirements
\end{itemize}

Changes may include installations, upgrades, patches, configuration changes, migrations, integrations, removals, and changes to service providers.

Routine administrative tasks that do not materially affect security, availability, performance, or functionality may follow an approved standard-change procedure.

\subsection{Change Classification}

Changes \should{} be classified according to their risk, urgency, complexity, and expected effect.

Change classifications may include:
\begin{itemize}
    \item \textbf{Standard change:} A low-risk, repeatable change with an approved procedure.
    \item \textbf{Normal change:} A planned change that requires review and approval before implementation.
    \item \textbf{Emergency change:} An urgent change required to address an active incident, critical vulnerability, service outage, or immediate operational risk.
\end{itemize}

The person implementing a change \must{} use the appropriate process for the change classification.

A change \must{} not be classified as standard merely to avoid review or approval requirements.

\subsection{Change Requests}

Normal changes \must{} be documented before implementation.

A change request \must{} include:
\begin{itemize}
    \item A description of the proposed change
    \item The reason and business purpose
    \item The affected systems, services, data, and users
    \item The expected security, availability, and operational effects
    \item The risks associated with the change
    \item The implementation steps
    \item The testing and validation plan
    \item The planned implementation date and maintenance window
    \item The rollback or recovery procedure
    \item The person responsible for implementation
    \item The required approvers
\end{itemize}

Changes \must{} be approved by an authorized person before implementation unless the emergency-change process applies.

\subsection{Change Assessment}

The change requester and system owner \must{} evaluate the potential effects of a proposed change.

The assessment \should{} consider:
\begin{itemize}
    \item Confidentiality, integrity, and availability
    \item Dependencies between systems and services
    \item Effects on users and business operations
    \item Security controls and monitoring
    \item Backup and recovery requirements
    \item Compatibility with existing hardware and software
    \item Data-protection and compliance requirements
    \item Required downtime or service interruption
    \item Risk of failure and the ability to reverse the change
\end{itemize}

High-risk changes \should{} receive additional technical review and management approval.

Changes affecting critical systems \should{} be scheduled during an approved maintenance window whenever reasonably practicable.

\subsection{Testing}

Changes \should{} be tested before production implementation.

Testing \should{} be performed in a development, test, staging, or other environment that is separate from production when reasonably practicable.

Testing \should{} verify:
\begin{itemize}
    \item The change performs its intended function
    \item Existing functionality remains available
    \item Security controls remain effective
    \item Data is not corrupted or lost
    \item Monitoring and logging continue to operate
    \item Backup and recovery procedures remain usable
    \item Dependent systems continue to function
\end{itemize}

Test results \must{} be documented for significant or high-risk changes.

Testing may be limited or skipped for an emergency change when delay would create greater risk. The reason \must{} be documented.

\subsection{Implementation}

Changes \must{} be implemented according to the approved change plan.

The person implementing the change \must{} verify that:
\begin{itemize}
    \item The correct system and environment are being changed
    \item A current backup or recovery point is available when appropriate
    \item Required personnel are available
    \item Required approvals have been obtained
    \item The approved implementation steps are followed
    \item The implementation is documented
\end{itemize}

Unplanned changes made during implementation \must{} be documented and reviewed.

Changes \must{} not be implemented using unauthorized accounts, unapproved software, or undocumented procedures.

\subsection{Rollback and Recovery}

Each significant change \must{} have a documented rollback or recovery procedure.

The rollback procedure \should{} identify:
\begin{itemize}
    \item Conditions that require rollback
    \item The person authorized to initiate rollback
    \item The steps required to restore the previous state
    \item Required backups or recovery points
    \item Expected service interruption
    \item Validation steps after rollback
    \item Communication requirements
\end{itemize}

If a change fails, causes an unacceptable security or operational issue, or cannot be validated, the change \should{} be rolled back when reasonably practicable.

\subsection{Post-Implementation Review}

Significant changes \must{} be reviewed after implementation.

The post-implementation review \should{} confirm:
\begin{itemize}
    \item The change achieved its intended purpose
    \item Systems and services are operating correctly
    \item Security controls remain effective
    \item Monitoring and logging are functioning
    \item No unauthorized changes were introduced
    \item Required documentation has been updated
    \item Any incidents or unexpected effects were addressed
\end{itemize}

Failed, partially completed, and emergency changes \must{} receive a post-implementation review.

The change record \must{} be updated with the implementation result, validation results, problems encountered, rollback actions, and final status.

\subsection{Emergency Changes}

Emergency changes may be implemented without normal advance approval when necessary to:
\begin{itemize}
    \item Contain or respond to a security incident
    \item Address an actively exploited vulnerability
    \item Restore a critical service
    \item Prevent immediate data loss
    \item Prevent significant harm to systems, users, or operations
\end{itemize}

Emergency changes \must{} be authorized by the responsible manager, system owner, incident leader, or other designated authority whenever reasonably practicable.

Emergency changes \must{} be documented as soon as possible and reviewed after implementation.

\subsection{Separation of Duties}

Where reasonably practicable, the person requesting a high-risk change \should{} not be the only person approving and implementing that change.

Changes affecting privileged access, security controls, financial systems, sensitive data, or critical services \should{} receive independent review.

No person \should{} approve their own high-risk change without additional review or documented justification.

\subsection{Change Records}

Change records \must{} be retained according to the organization's record-retention requirements.

Change records \should{} include:
\begin{itemize}
    \item Change identification number
    \item Request date
    \item Requester and responsible implementer
    \item Affected systems and services
    \item Change classification
    \item Risk assessment
    \item Approval records
    \item Planned and actual implementation dates
    \item Testing results
    \item Rollback or recovery information
    \item Implementation outcome
    \item Post-implementation review
\end{itemize}

System documentation, configuration baselines, network diagrams, operating procedures, and recovery documentation \must{} be updated after significant changes.

\subsection{Unauthorized Changes}

Unauthorized changes \must{} be investigated and documented.

The investigation \should{} determine:
\begin{itemize}
    \item What was changed
    \item When the change occurred
    \item Who or what performed the change
    \item Why the change occurred
    \item Whether security or availability was affected
    \item Whether rollback or remediation is required
    \item Whether additional controls are needed
\end{itemize}

Unauthorized changes that may indicate compromise, misuse, or loss of control \must{} be reported according to the incident-response procedures.